\chapter{Introduction}
\label{cap:intro}

Sensor deployments for environmental and infrastructure monitoring routinely
produce time-series data at a scale that makes naive range-scan query
execution impractical. A query asking for the average temperature across a
thousand sensors over the last month may need to scan billions of
individual readings if the underlying storage engine has no way to exploit
the temporal and sensor-wise structure of the data \cite{Jensen2017}.

\section{Motivation}
Existing time-series database systems typically rely on generic indexing
structures borrowed from relational systems, which do not exploit the fact
that sensor readings arrive at regular intervals and are naturally
partitioned by sensor identity and time window \cite{Pelkonen2015}.

\section{Contributions}
This dissertation makes the following contributions:
\begin{itemize}
  \item A query planning strategy that partitions scans along both the
    sensor and time dimensions, exploiting known sampling regularity to
    prune irrelevant partitions before execution.
  \item An implementation evaluated against three large-scale environmental
    sensing datasets with heterogeneous sampling rates.
  \item An analysis of when the proposed planner's gains are largest,
    showing the biggest improvements on queries aggregating across many
    sensors over long time windows.
\end{itemize}

\section{Dissertation Outline}
Chapter~\ref{cap:related} reviews prior work on time-series database
systems and query planning. Later chapters (omitted from this starter
document) would present the planning algorithm, implementation, and
experimental results in detail.
